\chapter{Introduction}
\label{ch:intro}

Recommender systems built on \textbf{matrix factorization} (MF) represent
users and items as vectors in a shared latent space and predict a
user-item affinity as their inner product. The technique is well
understood for \emph{explicit} feedback, where a rating directly measures
preference and the large mass of unrated pairs can be safely ignored during
training. Most real deployments, however, do not have ratings. They have
\emph{implicit} signals: clicks, purchases, views, reviews, which are only
ever positive. A user who never interacted with an item might dislike it,
might simply not know it exists, or might not have had the opportunity to
see it. This one-class problem, together with the sheer number of
unobserved pairs (millions of users times tens of thousands of items),
makes implicit-feedback MF both statistically and computationally harder
than its explicit counterpart.

\section{The paper}

This project is built around \citet{he2016eals}, \emph{Fast Matrix
Factorization for Online Recommendation with Implicit Feedback} (SIGIR
2016). The paper attacks the implicit-feedback problem on two fronts.
First, it argues that treating every unobserved entry as an equally weak
negative signal, the standard simplification introduced by
\citet{hu2008implicit}, throws away information a content provider usually
has for free: an item that was shown to many users but rarely interacted
with is a much stronger negative signal than an item nobody has seen. The
paper turns this intuition into a closed-form, popularity-based weighting
scheme for missing entries. Second, and this is the part this report
focuses on, it observes that once the missing data is weighted
non-uniformly, the classical ALS solver becomes prohibitively slow. That
solver updates a whole user or item vector at once by inverting a
$K\times K$ matrix, where $K$ is the number of latent factors each user and
item is represented with, and the memoization trick that makes
uniform-weight ALS tractable no longer applies once the weights vary per
item. The authors' fix, \textbf{eALS} (element-wise ALS), updates one
latent factor at a time in closed form. This avoids matrix inversion
entirely and reduces the per-iteration complexity from
$O((M+N)K^3 + |\mathcal{R}|K^2)$ to $O((M+N)K^2+|\mathcal{R}|K)$, where $M$
and $N$ are the numbers of users and items and $\mathcal{R}$ is the set of
observed interactions. Chapter~\ref{ch:formulation} restates all of this
notation formally before deriving eALS's update rule in
Chapter~\ref{ch:methods}.

\medskip
The paper states that this element-wise update is \emph{embarrassingly
parallel}: each user's (item's) update only depends on quantities that stay
frozen for the duration of a sweep. The authors' own implementation is
single-threaded Java, though, and no distributed version is provided.

\section{What this project does}

Following the assignment brief, this report:

\begin{enumerate}
  \item implements eALS (Algorithm~1 of the paper) as a genuine PySpark job
        rather than a single-machine simulation dressed up in Spark syntax:
        every update in every iteration is dispatched as an RDD
        transformation and executed by Spark's own task scheduler across
        partitions;
  \item evaluates it on the Amazon ``Movies and TV'' dataset, one of the two
        datasets used in the original paper (the other, the Yelp Challenge
        dataset, is no longer distributed by Yelp in the form the paper
        used, so this report does not attempt to reconstruct it);
  \item reproduces the paper's own experimental protocol as closely as
        practical (10-core filtering, leave-one-out evaluation, Hit Ratio
        and NDCG) and extends it with a dedicated scalability study
        against the number of latent factors, the degree of parallelism,
        and the training-set size, since this is the property the paper's
        central claim rests on but never measures directly in a
        distributed setting.
\end{enumerate}

\medskip
The paper's title promises \emph{online} recommendation, which in the
authors' terminology refers to a second contribution: an incremental
update rule (their Algorithm~2) that refreshes the model in $O(K^2)$ time
given a single new interaction, without retraining from scratch. This
report deliberately does not implement that part: the assignment brief asks
for ``an efficient matrix factorization ALS algorithm'' that ``can be
easily parallelized'', which describes Algorithm~1, not the online
extension built on top of it. This scope decision is flagged explicitly
here rather than left to surface as an omission later.

\medskip
The remainder of this report is organized as follows. Chapter~\ref{ch:formulation}
formalizes the optimization problem eALS solves. Chapter~\ref{ch:methods}
derives the element-wise update rule and describes how it was mapped onto
Spark's execution model. Chapter~\ref{ch:protocol} describes the dataset,
preprocessing, and evaluation protocol. Chapters~\ref{ch:results} and
\ref{ch:scalability} report the accuracy and scalability experiments.
Chapter~\ref{ch:discussion} discusses what worked, what did not, and where
this implementation departs from the paper. Chapter~\ref{ch:conclusion}
concludes.
